\documentclass[11pt,a4paper]{article}

% ---- Packages ----
\usepackage[utf8]{inputenc}
\usepackage[T1]{fontenc}
\usepackage{amsmath,amssymb,amsthm}
\usepackage{mathtools}
\usepackage[margin=1in,headheight=14pt]{geometry}
\usepackage{hyperref}
\usepackage{enumitem}
\usepackage{xcolor}
\usepackage{tcolorbox}
\usepackage{fancyhdr}
\usepackage{booktabs}

% ---- Hyperref ----
\hypersetup{
    colorlinks=true,
    linkcolor=red,
    citecolor=blue,
    urlcolor=blue
}

% ---- Header ----
\pagestyle{fancy}
\fancyhf{}
\fancyhead[L]{\textit{Lesson 10: Stochastic Approximation}}
\fancyhead[R]{\thepage}
\renewcommand{\headrulewidth}{0.4pt}

% ---- Theorem Environments ----
\theoremstyle{definition}
\newtheorem{definition}{Definition}[section]
\newtheorem{example}{Example}[section]
\newtheorem{exercise}{Exercise}[section]

\theoremstyle{plain}
\newtheorem{theorem}{Theorem}[section]
\newtheorem{lemma}[theorem]{Lemma}
\newtheorem{proposition}[theorem]{Proposition}
\newtheorem{corollary}[theorem]{Corollary}

\theoremstyle{remark}
\newtheorem{remark}{Remark}[section]

% ---- Colored Boxes ----
\tcbuselibrary{theorems,skins,breakable}

\newtcolorbox{keyresult}[1][]{
    colback=green!5!white,
    colframe=green!50!black,
    fonttitle=\bfseries,
    title=#1,
    breakable
}

\newtcolorbox{rlconnection}[1][]{
    colback=blue!5!white,
    colframe=blue!50!black,
    fonttitle=\bfseries,
    title=#1,
    breakable
}

\newtcolorbox{warning}[1][]{
    colback=red!5!white,
    colframe=red!50!black,
    fonttitle=\bfseries,
    title=#1,
    breakable
}

% ---- Operators ----
\newcommand{\R}{\mathbb{R}}
\newcommand{\N}{\mathbb{N}}
\newcommand{\Q}{\mathbb{Q}}
\newcommand{\Z}{\mathbb{Z}}
\newcommand{\C}{\mathbb{C}}
\newcommand{\Prob}{\mathbb{P}}
\newcommand{\E}{\mathbb{E}}
\newcommand{\Var}{\mathrm{Var}}
\newcommand{\indicator}{\mathbf{1}}
\newcommand{\cas}{\xrightarrow{\text{a.s.}}}
\newcommand{\cprob}{\xrightarrow{\Prob}}
\newcommand{\cdist}{\xrightarrow{d}}

% ---- Title ----
\title{Lesson 31: Robbins--Monro and the ODE Method}
\author{Curriculum: A Rigorous Learning Path to Multi-Agent Deep RL}
\date{}

\begin{document}
\maketitle
\tableofcontents
\newpage

\setcounter{section}{0}

\section{Introduction}
\label{sec:intro}

The central challenge in reinforcement learning is optimization under uncertainty: an
agent must find an optimal policy when the environment dynamics are unknown and the only
information available comes through noisy, sequential interactions. This is precisely the
setting of \emph{stochastic approximation} --- finding the root (or optimum) of a
function that can only be observed through noisy samples.

The field of stochastic approximation was born with the seminal 1951 paper of Robbins
and Monro, who proposed an iterative scheme for finding the root of a regression
function from noisy observations. This idea was extended to optimization by Kiefer and
Wolfowitz, and the theoretical foundations were deepened by Ljung, Kushner and Clark,
and Borkar, who established the powerful \emph{ODE method}: under appropriate conditions,
a stochastic approximation algorithm converges to the same limit as the solution of an
associated ordinary differential equation.

\medskip
\noindent\textbf{Why this matters for RL.} Consider the following RL update rules:
\begin{itemize}[leftmargin=2em]
\item \textbf{TD(0):} $V(s_n) \leftarrow V(s_n) + \alpha_n \bigl[R_{n+1} + \gamma V(s_{n+1}) - V(s_n)\bigr]$
\item \textbf{Q-learning:} $Q(s_n, a_n) \leftarrow Q(s_n,a_n) + \alpha_n \bigl[R_{n+1} + \gamma \max_{a'} Q(s_{n+1},a') - Q(s_n,a_n)\bigr]$
\item \textbf{Actor-Critic:} Critic updates value, actor updates policy, at \emph{different rates}
\item \textbf{REINFORCE:} $\theta_{n+1} = \theta_n + \alpha_n \widehat{\nabla J(\theta_n)}$
\end{itemize}
Every one of these is a stochastic approximation iteration. The convergence theorems
in this document therefore provide the theoretical foundation for understanding
\emph{when} and \emph{why} RL algorithms converge.

\medskip
\noindent\textbf{Document outline.}
Section~\ref{sec:robbins-monro} introduces the Robbins--Monro algorithm and proves its
convergence under classical conditions. Section~\ref{sec:ode-method} develops the ODE
method, the most powerful general-purpose tool for SA convergence analysis.
Section~\ref{sec:convergence-rate} establishes the convergence rate via the CLT for SA
and Polyak--Ruppert averaging. Section~\ref{sec:two-timescale} presents Borkar's
two-timescale SA theory, which is the backbone of actor-critic algorithms.
Section~\ref{sec:linear-sa} treats the important special case of linear SA.
Section~\ref{sec:rl-applications} maps the general theory to specific RL algorithms
in detail. Section~\ref{sec:exercises} provides exercises, and
Section~\ref{sec:summary} summarizes all key results.

\medskip
\noindent\textbf{Prerequisites.} Probability spaces, conditional expectation,
martingale difference sequences, the law of large numbers, and basic ODE theory
(existence and uniqueness of solutions, Lyapunov stability).

%% ============================================================
%% SECTION 2: THE ROBBINS-MONRO ALGORITHM
%% ============================================================

\section{The Robbins--Monro Algorithm}
\label{sec:robbins-monro}

\subsection{Problem formulation}

Suppose we wish to find a root of a function $h : \R^d \to \R^d$, i.e., a point
$x^* \in \R^d$ such that $h(x^*) = 0$. The difficulty is that we cannot evaluate $h$
directly. Instead, for each query point $x$, we observe a noisy sample
\[
  H(x, Y) \quad \text{where} \quad h(x) = \E[H(x, Y)]
\]
and $Y$ is a random variable (or random element) representing the noise source.

\begin{example}[Stochastic root-finding]
\label{ex:root-finding}
Let $h(\theta) = \E[f(X) - \theta]$ where $X$ is a random variable with unknown
distribution. Then $h(\theta^*) = 0$ when $\theta^* = \E[f(X)]$. Given i.i.d.\ samples
$X_1, X_2, \ldots$, the noisy observation is $H(\theta, X_n) = f(X_n) - \theta$.
\end{example}

\begin{definition}[Step-size sequence]
\label{def:step-size}
A \emph{step-size sequence} (or \emph{learning rate schedule}) is a sequence
$\{a_n\}_{n \geq 0}$ of positive real numbers satisfying:
\begin{enumerate}[label=(\roman*)]
\item $\displaystyle\sum_{n=0}^{\infty} a_n = \infty$ \quad (the steps are large enough to reach any target),
\item $\displaystyle\sum_{n=0}^{\infty} a_n^2 < \infty$ \quad (the steps are small enough that noise averages out).
\end{enumerate}
\end{definition}

\begin{remark}[Why both conditions are needed]
\label{rem:step-size-intuition}
Condition (i) ensures the iterates can travel an unbounded total distance --- without
it, the algorithm could ``freeze'' before reaching $x^*$ regardless of the initial
condition. Condition (ii) controls the accumulated noise variance: since the noise at
step $n$ is scaled by $a_n$, the total noise variance is proportional to
$\sum a_n^2 \Var(\text{noise})$, which must be finite for convergence. The canonical
choice is $a_n = c / n$ for some $c > 0$ (since $\sum 1/n = \infty$ while
$\sum 1/n^2 < \infty$). More generally, $a_n = c / n^\alpha$ works for any
$\alpha \in (1/2, 1]$.
\end{remark}

\subsection{The iteration}

\begin{definition}[Robbins--Monro iteration]
\label{def:rm-iteration}
Given an initial point $x_0 \in \R^d$, the \emph{Robbins--Monro iteration} is
\begin{equation}
\label{eq:rm}
  x_{n+1} = x_n + a_n H(x_n, Y_n), \quad n \geq 0,
\end{equation}
where $\{a_n\}$ satisfies Definition~\ref{def:step-size} and $Y_0, Y_1, \ldots$ are
the noise samples.
\end{definition}

We decompose the update into its ``signal'' and ``noise'' components. Define the
\emph{martingale difference noise}
\begin{equation}
\label{eq:mds-noise}
  M_{n+1} \coloneqq H(x_n, Y_n) - h(x_n).
\end{equation}
Then the iteration becomes
\begin{equation}
\label{eq:rm-decomposed}
  x_{n+1} = x_n + a_n\bigl[h(x_n) + M_{n+1}\bigr].
\end{equation}
This is the canonical form we shall use throughout: the iterate moves in the direction
$h(x_n)$ (the ``drift'') plus additive noise $M_{n+1}$.

\subsection{Classical convergence theorem}

We state the assumptions precisely and then prove convergence.

\begin{definition}[Standing assumptions for Robbins--Monro]
\label{def:rm-assumptions}
Let $\mathcal{F}_n = \sigma(x_0, Y_0, \ldots, Y_{n-1})$ be the natural filtration.
We assume:
\begin{enumerate}[label=(A\arabic*)]
\item \label{A1} \textbf{Step sizes:} $\{a_n\}$ satisfies Definition~\ref{def:step-size}.
\item \label{A2} \textbf{Martingale difference noise:} $\E[M_{n+1} \mid \mathcal{F}_n] = 0$ a.s.\ for all $n$.
\item \label{A3} \textbf{Bounded noise variance:} There exists $K > 0$ such that
  $\E[\|M_{n+1}\|^2 \mid \mathcal{F}_n] \leq K(1 + \|x_n\|^2)$ a.s.
\item \label{A4} \textbf{Stability (monotonicity):} There exists a unique root
  $x^* \in \R^d$ with $h(x^*) = 0$, and $\langle x - x^*, h(x) \rangle < 0$ for
  all $x \neq x^*$. (That is, $h$ points ``inward'' toward $x^*$.)
\item \label{A5} \textbf{Linear growth:} $\|h(x)\| \leq K(1 + \|x\|)$ for all $x$.
\end{enumerate}
\end{definition}

\begin{theorem}[Robbins--Monro convergence]
\label{thm:rm-convergence}
Under assumptions \ref{A1}--\ref{A5}, the Robbins--Monro iteration~\eqref{eq:rm}
satisfies $x_n \cas x^*$.
\end{theorem}

\begin{proof}
We use the Lyapunov function $V(x) = \|x - x^*\|^2$. Define $e_n = x_n - x^*$.
Then
\begin{align*}
  \|e_{n+1}\|^2
    &= \|e_n + a_n(h(x_n) + M_{n+1})\|^2 \\
    &= \|e_n\|^2 + 2a_n \langle e_n, h(x_n)\rangle
      + 2a_n \langle e_n, M_{n+1}\rangle \\
    &\quad + a_n^2 \|h(x_n) + M_{n+1}\|^2.
\end{align*}
Taking conditional expectation with respect to $\mathcal{F}_n$ and using \ref{A2}:
\begin{align}
  \E[\|e_{n+1}\|^2 \mid \mathcal{F}_n]
    &= \|e_n\|^2 + 2a_n \langle e_n, h(x_n)\rangle
      + a_n^2 \E[\|h(x_n) + M_{n+1}\|^2 \mid \mathcal{F}_n]. \label{eq:rm-lyap}
\end{align}
For the last term, by the inequality $\|a+b\|^2 \leq 2\|a\|^2 + 2\|b\|^2$:
\[
  \E[\|h(x_n) + M_{n+1}\|^2 \mid \mathcal{F}_n]
    \leq 2\|h(x_n)\|^2 + 2\E[\|M_{n+1}\|^2 \mid \mathcal{F}_n].
\]
Using \ref{A3} and \ref{A5}:
\[
  2\|h(x_n)\|^2 + 2\E[\|M_{n+1}\|^2 \mid \mathcal{F}_n]
    \leq 2K^2(1 + \|x_n\|^2) + 2K(1 + \|x_n\|^2)
    \leq C(1 + \|x_n\|^2)
\]
for some constant $C > 0$. Since $\|x_n\|^2 \leq 2\|e_n\|^2 + 2\|x^*\|^2$, we
obtain
\begin{equation}
\label{eq:rm-lyap2}
  \E[\|e_{n+1}\|^2 \mid \mathcal{F}_n]
    \leq \|e_n\|^2 + 2a_n \langle e_n, h(x_n)\rangle + a_n^2 C'(1 + \|e_n\|^2)
\end{equation}
for some $C' > 0$.

By \ref{A4}, $\langle e_n, h(x_n) \rangle < 0$ when $x_n \neq x^*$. Define
$V_n = \|e_n\|^2$ and
\[
  W_n = \sum_{k=0}^{n-1} (-2a_k) \langle e_k, h(x_k)\rangle, \quad
  U_n = \sum_{k=0}^{n-1} a_k^2 C'(1 + \|e_k\|^2).
\]
Then $V_n + W_n - U_n$ is a (nonnegative when adjusted) supermartingale. More
precisely, writing $\beta_n = -2\langle e_n, h(x_n)\rangle > 0$ (for $x_n \neq x^*$),
we have from~\eqref{eq:rm-lyap2}:
\[
  \E[V_{n+1} \mid \mathcal{F}_n] \leq V_n - a_n \beta_n + a_n^2 C'(1 + V_n + 2\|x^*\|^2).
\]

We apply the Robbins--Siegmund theorem (a standard supermartingale convergence result).

\medskip\noindent
\textbf{Robbins--Siegmund Theorem.} \emph{Let $\{V_n\}$, $\{\alpha_n\}$,
$\{\beta_n\}$, $\{u_n\}$ be nonneg\-a\-tive adapted sequences with
$\E[V_{n+1} \mid \mathcal{F}_n] \leq (1+\alpha_n)V_n - \beta_n + u_n$ and
$\sum \alpha_n < \infty$, $\sum u_n < \infty$ a.s. Then $V_n$ converges a.s.\ and
$\sum \beta_n < \infty$ a.s.}

\medskip
In our case, set $\alpha_n = a_n^2 C'$, $u_n = a_n^2 C'(1 + 2\|x^*\|^2)$, and
$\beta_n = a_n(-2\langle e_n, h(x_n)\rangle)$. Since $\sum a_n^2 < \infty$ by
\ref{A1}, we have $\sum \alpha_n < \infty$ and $\sum u_n < \infty$. By
Robbins--Siegmund, $V_n = \|e_n\|^2$ converges a.s.\ (to some finite limit) and
\begin{equation}
\label{eq:beta-summable}
  \sum_{n=0}^{\infty} a_n \bigl(-2\langle e_n, h(x_n)\rangle\bigr) < \infty
  \quad \text{a.s.}
\end{equation}

Suppose for contradiction that $V_n \to L > 0$ with positive probability on some
event $A$. On $A$, $\|e_n\| \to \sqrt{L} > 0$, so $x_n$ stays bounded away from
$x^*$. By the continuity of $h$ and \ref{A4}, there exists $\delta > 0$ such that
$-\langle e_n, h(x_n)\rangle \geq \delta$ for all large $n$ on $A$. Then
\[
  \sum_{n} a_n (-2\langle e_n, h(x_n)\rangle) \geq 2\delta \sum_n a_n = \infty,
\]
contradicting~\eqref{eq:beta-summable}. Therefore $V_n \to 0$ a.s., i.e.,
$x_n \cas x^*$.
\end{proof}

\begin{keyresult}[Robbins--Monro Convergence]
Under a monotonicity condition on $h$ (it points toward the root $x^*$),
martingale difference noise with bounded conditional variance, and the step-size
conditions $\sum a_n = \infty$, $\sum a_n^2 < \infty$, the iteration
$x_{n+1} = x_n + a_n[h(x_n) + M_{n+1}]$ converges almost surely to $x^*$.
\end{keyresult}

\begin{rlconnection}[SGD as Robbins--Monro]
Stochastic gradient descent (SGD) for minimizing $f(\theta) = \E_Z[\ell(\theta, Z)]$
is the Robbins--Monro iteration with $h(\theta) = -\nabla f(\theta)$ and
$H(\theta, Z) = -\nabla_\theta \ell(\theta, Z)$. The update is
$\theta_{n+1} = \theta_n - a_n \nabla_\theta \ell(\theta_n, Z_n)$. Under
convexity of $f$ (which implies the monotonicity condition \ref{A4} for $-\nabla f$),
Theorem~\ref{thm:rm-convergence} guarantees $\theta_n \cas \theta^*$.
\end{rlconnection}


%% ============================================================
%% SECTION 3: THE ODE METHOD
%% ============================================================

\section{The ODE Method}
\label{sec:ode-method}

The Robbins--Monro analysis requires the strong monotonicity condition \ref{A4},
which fails for many RL algorithms. The \emph{ODE method}, developed by Ljung (1977)
and Kushner--Clark (1978), replaces this with a more natural condition: convergence of
the SA iteration follows from the stability of the associated ODE. This is the
primary tool for analyzing RL convergence.

\subsection{The associated ODE}

\begin{definition}[Associated ODE]
\label{def:associated-ode}
Given the SA iteration~\eqref{eq:rm-decomposed}, the \emph{associated ODE} is
\begin{equation}
\label{eq:ode}
  \dot{x}(t) = h(x(t)), \quad x(0) = x_0.
\end{equation}
\end{definition}

The key insight is that for small step sizes, the SA iteration
$x_{n+1} \approx x_n + a_n h(x_n)$ behaves like an Euler discretization of the
ODE~\eqref{eq:ode} with time step $a_n$, plus noise that averages out.

\subsection{Interpolated trajectory}

To make the connection between the discrete SA iteration and the continuous ODE
precise, we construct an interpolated trajectory.

\begin{definition}[Interpolated trajectory]
\label{def:interpolated-traj}
Define the time instants $t_0 = 0$ and $t_n = \sum_{k=0}^{n-1} a_k$ for $n \geq 1$.
Note that $t_n \to \infty$ by condition (i) of Definition~\ref{def:step-size}. The
\emph{interpolated trajectory} $\bar{x} : [0, \infty) \to \R^d$ is the piecewise
constant function
\[
  \bar{x}(t) = x_n \quad \text{for } t \in [t_n, t_{n+1}).
\]
Equivalently, one may use the piecewise linear interpolation; both lead to the same
asymptotic analysis.
\end{definition}

\subsection{Assumptions for the ODE method}

\begin{definition}[ODE method assumptions]
\label{def:ode-assumptions}
We strengthen and modify the assumptions as follows:
\begin{enumerate}[label=(B\arabic*)]
\item \label{B1} \textbf{Step sizes:} As in \ref{A1}.
\item \label{B2} \textbf{Lipschitz drift:} $h : \R^d \to \R^d$ is Lipschitz
  continuous: $\|h(x) - h(y)\| \leq L\|x - y\|$ for all $x, y$ and some $L > 0$.
\item \label{B3} \textbf{Martingale difference noise:}
  $\E[M_{n+1} \mid \mathcal{F}_n] = 0$ a.s.\ and
  $\E[\|M_{n+1}\|^2 \mid \mathcal{F}_n] \leq K(1 + \|x_n\|^2)$ a.s.
\item \label{B4} \textbf{Stability (boundedness):} $\sup_n \|x_n\| < \infty$ a.s.
\item \label{B5} \textbf{Asymptotic stability of the ODE:} The ODE~\eqref{eq:ode}
  has a globally asymptotically stable equilibrium $x^*$ (i.e., $h(x^*) = 0$ and
  $x(t) \to x^*$ as $t \to \infty$ for every initial condition $x(0) \in \R^d$).
\end{enumerate}
\end{definition}

\begin{remark}[Stability assumption \ref{B4}]
\label{rem:stability}
Assumption \ref{B4} --- that the iterates remain bounded --- is the most delicate.
In practice, one either (a) proves it from properties of $h$ (see the Borkar--Meyn
theorem in Section~\ref{subsec:borkar-meyn}), (b) uses projection to a compact set,
or (c) verifies it empirically. For many RL algorithms, boundedness follows from the
structure of the MDP (e.g., bounded rewards and a discount factor $\gamma < 1$).
\end{remark}

\subsection{The noise vanishes asymptotically}

The first step in the ODE method is to show that the accumulated noise becomes
negligible.

\begin{lemma}[Noise term vanishes]
\label{lem:noise-vanishes}
Under \ref{B1}, \ref{B3}, and \ref{B4}, the process
$Z_n = \sum_{k=0}^{n-1} a_k M_{k+1}$ converges a.s. In particular, for any
$T > 0$ and $\varepsilon > 0$,
\[
  \sup_{t_n \leq s \leq t_n + T}
    \Bigl\|\sum_{k=n}^{m(s)-1} a_k M_{k+1}\Bigr\| \to 0 \quad \text{a.s.\ as } n \to \infty,
\]
where $m(s) = \min\{j : t_j > s\}$.
\end{lemma}

\begin{proof}
The sequence $Z_n = \sum_{k=0}^{n-1} a_k M_{k+1}$ is a martingale with respect to
$\mathcal{F}_n$. Its quadratic variation satisfies
\[
  \E[\|Z_n\|^2] = \sum_{k=0}^{n-1} a_k^2 \E[\|M_{k+1}\|^2]
    \leq \sum_{k=0}^{n-1} a_k^2 K(1 + \E[\|x_k\|^2]).
\]
By \ref{B4}, $\|x_k\|$ is bounded a.s., so $\E[\|x_k\|^2] \leq C$ for some $C > 0$.
Thus $\E[\|Z_n\|^2] \leq K(1 + C) \sum_{k=0}^{\infty} a_k^2 < \infty$ by \ref{B1}.
By the martingale convergence theorem, $Z_n$ converges a.s.

For the uniform convergence over windows of length $T$, note that for
$t_n \leq s \leq t_n + T$, the ``tail sum'' $\sum_{k=n}^{m(s)-1} a_k M_{k+1}$
is a portion of the convergent series $\sum a_k M_{k+1}$, so it tends to zero
uniformly over bounded time windows. More precisely, since $Z_n$ converges a.s.,
$\|Z_m - Z_n\| \to 0$ as $n, m \to \infty$, and the number of indices $k$ with
$t_n \leq t_k \leq t_n + T$ increases but $a_k \to 0$, giving the claimed uniform
convergence.
\end{proof}

\subsection{Main convergence theorem}

\begin{theorem}[Convergence via the ODE method]
\label{thm:ode-convergence}
Under assumptions \ref{B1}--\ref{B5}, the SA iteration
$x_{n+1} = x_n + a_n[h(x_n) + M_{n+1}]$ satisfies $x_n \cas x^*$.
\end{theorem}

\begin{proof}
The proof proceeds in three steps.

\medskip
\noindent\textbf{Step 1: Shifted interpolated trajectories.}
For each $n \geq 0$, define the shifted interpolated trajectory
$x^{(n)} : [0, \infty) \to \R^d$ by
\[
  x^{(n)}(0) = x_n, \quad
  x^{(n)}(t) = x_{n+k} \text{ for } t \in [t_{n+k} - t_n, t_{n+k+1} - t_n),
  \quad k \geq 0.
\]
This is the piecewise constant interpolation of the SA iterates starting from time
$t_n$, with the time axis reset to zero.

\medskip
\noindent\textbf{Step 2: Equicontinuity and convergence to ODE solutions.}
By \ref{B4}, the iterates $\{x_n\}$ are bounded a.s. The drift $h$ is Lipschitz
by \ref{B2}, so $\|h(x_n)\| \leq L\|x_n\| + \|h(0)\|$ is bounded. For any
$T > 0$ and $s, t \in [0, T]$ with $s < t$, the trajectory satisfies
\[
  x^{(n)}(t) - x^{(n)}(s) = \sum_{k : t_k \in [s+t_n, t+t_n)} a_k [h(x_k) + M_{k+1}].
\]
The ``signal'' part is bounded by $\|h\|_\infty \cdot (t-s)$ (since
$\sum a_k \approx t - s$ over the corresponding indices). By
Lemma~\ref{lem:noise-vanishes}, the noise part vanishes as $n \to \infty$. Therefore,
the family $\{x^{(n)}\}$ is equicontinuous on $[0, T]$ a.s.\ for large $n$.

By the Arzela--Ascoli theorem, every subsequence of $\{x^{(n)}\}$ has a further
subsequence that converges uniformly on $[0, T]$. Let $x^{(\infty)}$ be any such
limit. We claim $x^{(\infty)}$ satisfies the ODE~\eqref{eq:ode}.

Indeed, the SA iteration gives
\[
  x^{(n)}(t) = x^{(n)}(0) + \int_0^t h(x^{(n)}(s))\,ds + \varepsilon_n(t),
\]
where $\varepsilon_n(t)$ accounts for the discretization error and the noise, both
of which vanish as $n \to \infty$ (the discretization error vanishes because
$a_n \to 0$ by \ref{B1}, and the noise vanishes by Lemma~\ref{lem:noise-vanishes}).
Taking the limit and using the Lipschitz continuity of $h$ to pass through the
integral:
\[
  x^{(\infty)}(t) = x^{(\infty)}(0) + \int_0^t h(x^{(\infty)}(s))\,ds.
\]
Thus $x^{(\infty)}$ is a solution of the ODE~\eqref{eq:ode}.

\medskip
\noindent\textbf{Step 3: Convergence to the equilibrium.}
By \ref{B5}, the equilibrium $x^*$ is globally asymptotically stable for the ODE.
Therefore, for any solution $x(t)$ of the ODE, $x(t) \to x^*$ as $t \to \infty$.

Now we argue that $x_n \to x^*$ a.s. The set of limit points of $\{x_n\}$ is
nonempty (by boundedness \ref{B4}) and compact. From Step~2, any limit point of the
discrete sequence is also an initial condition for an ODE trajectory that must
converge to $x^*$. Moreover, the set of limit points is connected (as the image of
a connected set under a continuous map) and invariant under the ODE flow.

The only compact, connected, ODE-flow-invariant set that is attracted to $x^*$ under
a globally asymptotically stable system is $\{x^*\}$ itself. To see this, suppose
$\bar{x} \neq x^*$ is a limit point. Then there exists a subsequence $x_{n_k} \to
\bar{x}$. The shifted trajectory $x^{(n_k)}$ converges to an ODE solution starting
at $\bar{x}$, which satisfies $x(T) \to x^*$ for large $T$. But then the iterates
after $n_k$ must eventually be close to $x^*$, and by the same argument, they
cannot leave a neighborhood of $x^*$ (by Lyapunov stability of $x^*$). Hence $x^*$
is the unique limit point, and $x_n \cas x^*$.
\end{proof}

\begin{keyresult}[ODE Method -- Main Principle]
A stochastic approximation iteration $x_{n+1} = x_n + a_n[h(x_n) + M_{n+1}]$
with decreasing step sizes converges to the set of asymptotically stable equilibria
of the ODE $\dot{x} = h(x)$, provided the iterates remain bounded and the noise is
a martingale difference sequence.
\end{keyresult}

\subsection{The Borkar--Meyn stability theorem}
\label{subsec:borkar-meyn}

Assumption \ref{B4} (boundedness of iterates) is often the hardest to verify. The
Borkar--Meyn theorem provides sufficient conditions.

\begin{definition}[Borkar--Meyn assumptions]
\label{def:borkar-meyn-assumptions}
In addition to \ref{B1}--\ref{B3}, assume:
\begin{enumerate}[label=(C\arabic*)]
\item \label{C1} $h$ is Lipschitz and the ``scaled'' function
  $h_\infty(x) \coloneqq \lim_{r \to \infty} h(rx)/r$ exists uniformly on compact
  sets.
\item \label{C2} The ODE $\dot{x} = h_\infty(x)$ has the origin as its unique
  globally asymptotically stable equilibrium.
\end{enumerate}
\end{definition}

\begin{theorem}[Borkar--Meyn stability]
\label{thm:borkar-meyn}
Under \ref{B1}--\ref{B3} and \ref{C1}--\ref{C2}, the iterates $\{x_n\}$ are bounded
almost surely: $\sup_n \|x_n\| < \infty$ a.s.
\end{theorem}

\begin{proof}
We give the key steps of the proof.

\medskip
\noindent\textbf{Step 1: Rescaling.}
Suppose, for contradiction, that $\|x_n\| \to \infty$ along a subsequence with
positive probability. Let $r_n = \max(\|x_0\|, \|x_1\|, \ldots, \|x_n\|)$ and
consider the rescaled iterates $\hat{x}_n = x_n / r_n$.

\medskip
\noindent\textbf{Step 2: Rescaled iteration.}
Since $x_{n+1} = x_n + a_n[h(x_n) + M_{n+1}]$, dividing by $r_{n+1}$:
\[
  \hat{x}_{n+1} = \frac{r_n}{r_{n+1}} \hat{x}_n
    + \frac{a_n}{r_{n+1}} [h(r_n \hat{x}_n) + M_{n+1}].
\]
When $r_n \to \infty$, $h(r_n \hat{x}_n) / r_n \to h_\infty(\hat{x}_n)$ by \ref{C1},
and $r_n / r_{n+1} \to 1$ since $a_n \to 0$. The noise term $a_n M_{n+1}/r_{n+1}$
becomes negligible.

\medskip
\noindent\textbf{Step 3: Limiting ODE.}
By the same interpolation argument as in Theorem~\ref{thm:ode-convergence}, the
rescaled trajectory converges (along subsequences) to solutions of the ODE
$\dot{x} = h_\infty(x)$. But $\|\hat{x}_n\| \leq 1$ and $\|\hat{x}_n\| = 1$
infinitely often (when $r_n$ increases). So the limiting trajectory is confined to
the unit ball and visits its boundary.

\medskip
\noindent\textbf{Step 4: Contradiction.}
By \ref{C2}, every solution of $\dot{x} = h_\infty(x)$ converges to the origin.
But the limiting trajectory of $\hat{x}_n$ must repeatedly visit $\|\hat{x}\| = 1$,
contradicting the global asymptotic stability of the origin for the $h_\infty$ ODE.
Therefore $\sup_n \|x_n\| < \infty$ a.s.
\end{proof}

\begin{remark}[Practical significance]
\label{rem:borkar-meyn-practical}
The Borkar--Meyn theorem is crucial because it converts the verification of iterate
boundedness into a condition on the ``behavior of $h$ at infinity.'' For linear SA
where $h(x) = Ax + b$, we have $h_\infty(x) = Ax$, so \ref{C2} reduces to
requiring that $A$ is Hurwitz (all eigenvalues have negative real parts). This is
exactly the condition for TD(0) convergence with linear function approximation.
\end{remark}


%% ============================================================
%% SECTION 4: CONVERGENCE RATE ANALYSIS
%% ============================================================


\end{document}